\chapter{Pseudorandom Functions and CPA Security}\label{ch:prf-cpa}

At the end of the previous chapter we beat Shannon's bound: with a PRG, a short key encrypts a long message. But the scheme $\Enc(k, m) = G(k) \oplus m$ is still \textbf{one-time}: the pad $G(k)$ is always the same, so a key can protect a single message. This chapter removes that limitation:
\begin{itemize}
  \item \textbf{a stronger definition} $\to$ security when many messages are encrypted with the same key, even against an adversary who can obtain encryptions of messages of its choice (CPA security);
  \item \textbf{a new tool} $\to$ pseudorandom functions, a ``secret giant table'' that looks like a random function;
  \item \textbf{a scheme and its proof} $\to$ $\Enc(k, m) = (r, F_k(r) \oplus m)$ is CPA-secure;
  \item \textbf{a construction} $\to$ pseudorandom functions can be built from PRGs (GGM).
\end{itemize}

% ============================================================
%  Security against chosen-plaintext attacks
% ============================================================

\section{Security against chosen-plaintext attacks}

\subsection{The threat model}

In practice Alice and Bob agree on \textbf{one} key and use it for the whole conversation. Moreover, the adversary can often \textbf{make Alice encrypt messages of its choice}:
\begin{itemize}
  \item \textbf{e-mail} $\to$ Eve sends an e-mail to Alice, and Alice's server stores or forwards it encrypted with Alice's key;
  \item \textbf{web} $\to$ a malicious page makes the victim's browser send requests chosen by the attacker, which leave encrypted with the victim's key (attacks on TLS such as BEAST relied on this);
  \item \textbf{history} $\to$ during the Second World War, the British laid mines in chosen places, knowing that German reports would encrypt those place names with Enigma.
\end{itemize}
This is a \textbf{chosen-plaintext attack} (CPA). The adversary chooses the \textbf{messages}, never the key nor the randomness of the encryption, and sees only the ciphertexts.

As usual, giving the adversary more power gives a stronger guarantee: a scheme that resists an adversary who can encrypt whatever it wants also resists one who cannot.

\subsection{The definition}

\begin{frameddefn}[label={def:cpa-game}]{The CPA game $\mathrm{GAME}^{\text{cpa}}_{\Pi,\Adv}(\lambda, b)$}
  Let $b \in \{0,1\}$ be a fixed bit. The adversary $\Adv$ plays against a challenger $C$:
  \begin{enumerate}
    \item $C$ chooses a key $k \sample \Kspace$, \textbf{once};
    \item \textbf{encryption queries:} $\Adv$ sends messages $m$ of its choice and receives $\Enc(k, m)$, always with the \textbf{same} key $k$;
    \item \textbf{challenge:} $\Adv$ sends two messages $m_0, m_1$ of the same length and receives $c^* \leftarrow \Enc(k, m_b)$;
    \item $\Adv$ may ask more encryption queries, and finally outputs a bit $b'$, which is the output of the game.
  \end{enumerate}
  $\Adv$ is PPT, so it asks at most a polynomial number of queries.
\end{frameddefn}

\begin{center}
  \begin{tikzpicture}[x=1cm, y=1cm]
    \node (A) at (0,0) {$\Adv$};
    \node (C) at (7,0) {$C$};
    \node[lbl, anchor=west] at (7.3,-0.5) {$k \sample \Kspace$};
    \draw[msg] (0,-1.0) -- node[above] {$m$} (7,-1.0);
    \draw[msg] (7,-1.6) -- node[above] {$\Enc(k, m)$} (0,-1.6);
    \node[lbl] at (3.5,-2.0) {$\vdots$ \ (polynomially many queries)};
    \draw[msg] (0,-2.7) -- node[above] {$m_0, m_1$} (7,-2.7);
    \draw[msg] (7,-3.3) -- node[above] {$c^* \leftarrow \Enc(k, m_b)$} (0,-3.3);
    \node[lbl] at (3.5,-3.7) {$\vdots$ \ (more queries)};
    \draw[msg] (0,-4.4) -- node[above] {$b'$} (7,-4.4);
  \end{tikzpicture}
\end{center}

\begin{frameddefn}[label={def:cpa-security}]{CPA security}
  A symmetric-key encryption scheme $\Pi = (\Enc, \Dec)$ is \textbf{CPA-secure} if
  \[
    \mathrm{GAME}^{\text{cpa}}_{\Pi,\Adv}(\lambda, 0) \cind \mathrm{GAME}^{\text{cpa}}_{\Pi,\Adv}(\lambda, 1)
  \]
  that is, for every PPT adversary $\Adv$:
  \[
    \Big|\Pr\big[\mathrm{GAME}^{\text{cpa}}_{\Pi,\Adv}(\lambda, 0) = 1\big] - \Pr\big[\mathrm{GAME}^{\text{cpa}}_{\Pi,\Adv}(\lambda, 1) = 1\big]\Big| \leq \negl(\lambda)
  \]
\end{frameddefn}

It is the one-time game of \Cref{def:one-time-security} with one difference: the adversary has access to an \textbf{encryption oracle} under the same key, before and after the challenge.

\subsection{CPA security requires randomized encryption}

\begin{framedobs}[label={obs:deterministic-cpa}]{No deterministic scheme is CPA-secure}
  If $\Enc$ is deterministic, the same message under the same key always gives the same ciphertext. Then the adversary:
  \begin{enumerate}
    \item asks the encryption of $m_0$ and receives $c_0 = \Enc(k, m_0)$;
    \item sends $m_0, m_1$ (with $m_0 \neq m_1$) as challenge and receives $c^*$;
    \item outputs $b' = 0$ if $c^* = c_0$, and $b' = 1$ otherwise.
  \end{enumerate}
  It always wins: its advantage is $1$. So a CPA-secure scheme must be \textbf{randomized}: it needs \textbf{fresh randomness at every encryption}, so that the same message has many possible ciphertexts.
\end{framedobs}

In particular the scheme $G(k) \oplus m$ of \Cref{def:prg-otp} is not CPA-secure. We need a new tool.

% ============================================================
%  Pseudorandom functions
% ============================================================

\section{Pseudorandom functions}

\subsection{The idea: a secret giant table}

A PRG gives \textbf{one} long string that looks random. A pseudorandom function gives a whole \textbf{function} that looks random: we can ask for its value on \textbf{any} input, whenever we need it.

Think of a truly random function $R : \{0,1\}^n \to \{0,1\}^\ell$ as a \textbf{table}, with one row for each of the $2^n$ inputs and a uniformly random $\ell$-bit value in each row. Such a table cannot be stored: with $n = 128$ it has $2^{128}$ rows. A PRF \textbf{simulates} it with a short key: the value of row $x$ is computed on demand as $F_k(x)$.

\begin{frameddefn}[label={def:prf}]{Pseudorandom function}
  Let $\mathcal{F} = \{F_k : \{0,1\}^n \to \{0,1\}^\ell\}_{k \in \{0,1\}^\lambda}$ be a family of functions computable in polynomial time. $\mathcal{F}$ is a \textbf{pseudorandom function} (PRF) if, for every PPT adversary $\Adv$,
  \[
    \mathrm{REAL}_{\mathcal{F},\Adv}(\lambda) \cind \mathrm{RAND}_{\mathcal{F},\Adv}(\lambda)
  \]
  where in both games $\Adv$ asks queries $x \in \{0,1\}^n$ (polynomially many, each one possibly chosen depending on the previous answers) and finally outputs a bit:
  \begin{itemize}
    \item \textbf{$\mathrm{REAL}$} $\to$ the challenger chooses $k \sample \{0,1\}^\lambda$ once, and answers every query $x$ with $F_k(x)$;
    \item \textbf{$\mathrm{RAND}$} $\to$ the challenger answers every query $x$ with $R(x)$, where $R$ is a truly random function from $\{0,1\}^n$ to $\{0,1\}^\ell$.
  \end{itemize}
\end{frameddefn}

Three lengths appear in the definition:
\begin{itemize}
  \item \textbf{$\lambda$} $\to$ the length of the key $k$;
  \item \textbf{$n$} $\to$ the length of the input $x$;
  \item \textbf{$\ell$} $\to$ the length of the output $F_k(x)$.
\end{itemize}

\paragraph{How the challenger of RAND samples a random function.} It cannot write down the whole table, so it builds it \textbf{lazily}, one row at a time, keeping a list of the rows already filled:
\begin{itemize}
  \item if $x$ is a \textbf{new} input, it chooses $y \sample \{0,1\}^\ell$, writes $(x, y)$ in the list and answers $y$;
  \item if $x$ was \textbf{already asked}, it answers with the same $y$ as before.
\end{itemize}
The answers are distributed exactly as the values of a truly random function. Note that $R$ is a \textbf{function}: the same input always gets the same answer.

\begin{framedobs}{How many functions are there?}
  A table with $2^n$ rows, each containing one of $2^\ell$ values, can be filled in
  \[
    \underbrace{2^\ell \cdot 2^\ell \cdots 2^\ell}_{2^n \text{ rows}} = 2^{\ell \cdot 2^n}
  \]
  ways. So choosing a truly random function requires $\ell \cdot 2^n$ random bits, the whole table. A PRF chooses one function among only $2^\lambda$, with a key of $\lambda$ bits: a tiny fraction of all functions, which no efficient adversary can tell apart from a random one.
\end{framedobs}

\begin{framedobs}{PRFs are deterministic}
  For a fixed key, $F_k$ is a deterministic function: the same input always gives the same output. This is essential: in an encryption scheme, Bob must recompute exactly the value computed by Alice. The randomness of the encryption will come from \textbf{where} the PRF is evaluated, not from the PRF itself.
\end{framedobs}

\paragraph{Comparison with pairwise independent hashing.} Both are families of functions indexed by a random key, but the guarantees are very different:

\begin{table}[H]
  \centering
  \small
  \begin{tabularx}{\textwidth}{@{}lXX@{}}
    \toprule
    & \textbf{Pairwise independent family} & \textbf{PRF} \\
    \midrule
    Guarantee & uniform outputs on \textbf{two} inputs & looks random on \textbf{polynomially many} inputs \\
    Against & unbounded adversaries & efficient adversaries \\
    Queries & fixed in advance & adaptive \\
    Key & grows with the number of uses & short, reused \\
    \bottomrule
  \end{tabularx}
\end{table}

In practice, block ciphers such as AES are used (and assumed to behave) as PRFs.

% ============================================================
%  CPA-secure encryption from a PRF
% ============================================================

\section{CPA-secure encryption from a PRF}

\subsection{The scheme}

The idea: keep the One-Time Pad structure, but take a \textbf{different pad for every message}, by evaluating the PRF on a \textbf{fresh random point}.

\begin{frameddefn}[label={def:prf-enc}]{Encryption from a PRF}
  Let $\mathcal{F} = \{F_k : \{0,1\}^n \to \{0,1\}^\ell\}_{k \in \{0,1\}^\lambda}$ be a PRF. The scheme $\Pi = (\Enc, \Dec)$ has key space $\{0,1\}^\lambda$, message space $\Mspace = \{0,1\}^\ell$ and ciphertext space $\{0,1\}^n \times \{0,1\}^\ell$:
  \[
    \begin{aligned}
      \Enc(k, m) &= (c_1, c_2) = \big(r,\ F_k(r) \oplus m\big) && \text{with } r \sample \{0,1\}^n \text{ fresh at every encryption} \\
      \Dec\big(k, (c_1, c_2)\big) &= F_k(c_1) \oplus c_2
    \end{aligned}
  \]
\end{frameddefn}

\begin{itemize}
  \item \textbf{Correctness} $\to$ $F_k(c_1) \oplus c_2 = F_k(r) \oplus F_k(r) \oplus m = m$.
  \item \textbf{$r$ is the label, $F_k(r)$ is the pad} $\to$ $r$ travels in clear and tells Bob which pad to recompute; the pad itself cannot be computed without $k$.
  \item \textbf{$r$ can be public} $\to$ in the PRF game the adversary even \textbf{chooses} the inputs, and still $F_k(x)$ looks random to it. The secret is $k$, not $r$ (as with the public seed of an extractor).
  \item \textbf{The scheme is randomized} $\to$ the randomness comes from $r$; the PRF is deterministic. The same message encrypted twice gives two different ciphertexts.
  \item \textbf{Lengths} $\to$ the message has $\ell$ bits because it is XORed with the pad; the ciphertext has $n + \ell$ bits, $n$ more than the message: the price of randomization.
\end{itemize}

\begin{framedex}{Encrypting the same message twice}
  Alice encrypts the same message $m$ twice with the same key:
  \begin{itemize}
    \item first time: $r = 7$, ciphertext $\big(7,\ F_k(7) \oplus m\big)$;
    \item second time: $r = 42$, ciphertext $\big(42,\ F_k(42) \oplus m\big)$.
  \end{itemize}
  The two pads $F_k(7)$ and $F_k(42)$ are different and look independent, so Eve cannot tell that the message is the same. Bob reads the labels $7$ and $42$, recomputes the two pads with $k$ and removes them.
\end{framedex}

\subsection{Security}

\begin{framedthm}[label={thm:prf-cpa}]{Encryption from a PRF is CPA-secure}
  If $\mathcal{F}$ is a PRF, then the scheme of \Cref{def:prf-enc} is CPA-secure.
\end{framedthm}

The proof is a game hopping with three kinds of games. Let $q = \poly(\lambda)$ be the total number of ciphertexts the adversary sees (the answers to its queries and the challenge), with labels $r_1, \ldots, r_q$.

\begin{table}[H]
  \centering
  \small
  \begin{tabularx}{\textwidth}{@{}lXl@{}}
    \toprule
    \textbf{Game} & \textbf{Every ciphertext (queries and challenge)} & \textbf{Hop} \\
    \midrule
    $\mathrm{GAME}(\lambda, b)$ & $(r,\ F_k(r) \oplus m)$, with $k \sample \{0,1\}^\lambda$ & \\
    $\mathrm{HYB}(\lambda, b)$ & $(r,\ R(r) \oplus m)$, with $R$ a truly random function & $\cind$ (PRF security) \\
    $\mathrm{HYB}'(\lambda)$ & $(r, u)$, with $r \sample \{0,1\}^n$ and $u \sample \{0,1\}^\ell$: pure noise & $\approx_s$ (collisions are rare) \\
    \bottomrule
  \end{tabularx}
  \caption{The games of the proof. $\mathrm{HYB}'$ does not depend on $b$: the challenge is random whatever message was chosen.}
\end{table}

\paragraph{The plan.} $\mathrm{HYB}'$ does not depend on $b$, so it is the common middle point:
\[
  \mathrm{GAME}(\lambda, 0) \;\cind\; \mathrm{HYB}(\lambda, 0) \;\approx_s\; \mathrm{HYB}'(\lambda) \;\approx_s\; \mathrm{HYB}(\lambda, 1) \;\cind\; \mathrm{GAME}(\lambda, 1)
\]

\subsubsection{Step 1: from the PRF to a random function}

\begin{framedlem}[label={lem:game-hyb}]{}
  For $b \in \{0,1\}$: $\mathrm{GAME}(\lambda, b) \cind \mathrm{HYB}(\lambda, b)$.
\end{framedlem}

\begin{proof}
  Suppose a PPT $\Adv$ tells $\mathrm{GAME}(\lambda, b)$ and $\mathrm{HYB}(\lambda, b)$ apart with non-negligible advantage. We build a PPT $B$ that plays the PRF game: $B$ has access to an oracle $O$, which is either $F_k$ (game $\mathrm{REAL}$) or a random function $R$ (game $\mathrm{RAND}$), and simulates the CPA challenger for $\Adv$.

  \begin{center}
    \begin{tikzpicture}[x=1cm, y=1cm]
      \node (A) at (0,0) {$\Adv$};
      \node (B) at (5,0) {$B$};
      \node (O) at (10,0) {$O \in \{F_k, R\}$};
      \draw[msg] (0,-0.7) -- node[above] {$m$} (5,-0.7);
      \node[lbl] at (5,-1.05) {$r \sample \{0,1\}^n$};
      \draw[msg] (5,-1.5) -- node[above] {$r$} (10,-1.5);
      \draw[msg] (10,-2.1) -- node[above] {$y = O(r)$} (5,-2.1);
      \draw[msg] (5,-2.7) -- node[above] {$(r,\ y \oplus m)$} (0,-2.7);
      \draw[msg] (0,-3.4) -- node[above] {$m_0, m_1$} (5,-3.4);
      \draw[msg] (5,-4.0) -- node[above] {$(r^*,\ O(r^*) \oplus m_b)$} (0,-4.0);
      \draw[msg] (0,-4.7) -- node[above] {$b'$} (5,-4.7);
      \draw[msg] (5,-4.7) -- node[above] {$b'$} (10,-4.7);
    \end{tikzpicture}
  \end{center}

  For every encryption query $m$ (and for the challenge, with $m_b$), $B$ chooses a fresh $r \sample \{0,1\}^n$, asks $O(r)$ and answers $(r, O(r) \oplus m)$. At the end it outputs the bit of $\Adv$.
  \begin{itemize}
    \item If $O = F_k$, $\Adv$ is playing \textbf{exactly} $\mathrm{GAME}(\lambda, b)$.
    \item If $O = R$, $\Adv$ is playing \textbf{exactly} $\mathrm{HYB}(\lambda, b)$.
  \end{itemize}
  The simulation is perfect, so the advantage of $B$ in distinguishing $\mathrm{REAL}$ from $\mathrm{RAND}$ equals the advantage of $\Adv$, which is non-negligible. $B$ is PPT (one oracle query and one XOR per ciphertext). This contradicts the security of the PRF.
\end{proof}

\subsubsection{Step 2: from a random function to pure noise}

In $\mathrm{HYB}$ the pad of each ciphertext is $R(r_i)$. If all the labels are distinct, each $R(r_i)$ is a \textbf{new} row of the random table: a fresh uniform value, independent of everything else. Then $R(r_i) \oplus m_i$ is uniform whatever $m_i$ is (a One-Time Pad with a fresh pad), and each ciphertext is pure noise, exactly as in $\mathrm{HYB}'$. The two games differ only if two labels collide, because then two ciphertexts share the same pad (and XORing them gives $m_i \oplus m_j$). Call this event
\[
  \mathrm{BAD} = \big[\exists\, i \neq j : r_i = r_j\big]
\]
The labels are uniform and independent in both games, so $\Pr[\mathrm{BAD}]$ is the same in $\mathrm{HYB}$ and in $\mathrm{HYB}'$.

\begin{framedlem}[label={lem:difference}]{Difference lemma}
  Let $H$ and $H'$ be two games and $\mathrm{BAD}$ an event such that $\Pr[\mathrm{BAD}]$ is the same in both, and, conditioned on $\overline{\mathrm{BAD}}$, the view of the adversary has the same distribution in both. Then
  \[
    \big|\Pr[H = 1] - \Pr[H' = 1]\big| \leq \Pr[\mathrm{BAD}]
  \]
\end{framedlem}

\begin{proof}
  \leavevmode
  \begin{enumerate}
    \item \textbf{Split on $\mathrm{BAD}$} (law of total probability), in both games:
    \[
      \Pr[H = 1] = \Pr[H = 1 \wedge \mathrm{BAD}] + \Pr[H = 1 \wedge \overline{\mathrm{BAD}}]
    \]
    and the same for $H'$.
    \item \textbf{Regroup and use the triangle inequality:}
    \[
      \begin{aligned}
        \big|\Pr[H = 1] - \Pr[H' = 1]\big| &\leq
        \underbrace{\big|\Pr[H = 1 \wedge \mathrm{BAD}] - \Pr[H' = 1 \wedge \mathrm{BAD}]\big|}_{|A|} \\
        &\quad + \underbrace{\big|\Pr[H = 1 \wedge \overline{\mathrm{BAD}}] - \Pr[H' = 1 \wedge \overline{\mathrm{BAD}}]\big|}_{|B|}
      \end{aligned}
    \]
    \item \textbf{The part without $\mathrm{BAD}$ vanishes.} With conditional probabilities,
    \[
      B = \Pr[\overline{\mathrm{BAD}}] \cdot \Big(\Pr[H = 1 \mid \overline{\mathrm{BAD}}] - \Pr[H' = 1 \mid \overline{\mathrm{BAD}}]\Big) = 0
    \]
    because, conditioned on $\overline{\mathrm{BAD}}$, the adversary sees the same distribution in both games, and an algorithm that receives the same distribution outputs $1$ with the same probability.
    \item \textbf{The part with $\mathrm{BAD}$ is at most $\Pr[\mathrm{BAD}]$.} We know nothing about the adversary when $\mathrm{BAD}$ happens, but a difference of two probabilities is at most $1$:
    \[
      |A| = \Pr[\mathrm{BAD}] \cdot \Big|\Pr[H = 1 \mid \mathrm{BAD}] - \Pr[H' = 1 \mid \mathrm{BAD}]\Big| \leq \Pr[\mathrm{BAD}] \qedhere
    \]
  \end{enumerate}
\end{proof}

In words: if two games are identical until a bad event happens, their distance is at most the probability of the bad event.

\begin{framedlem}[label={lem:hyb-hybprime}]{}
  For $b \in \{0,1\}$: $\big|\Pr[\mathrm{HYB}(\lambda, b) = 1] - \Pr[\mathrm{HYB}'(\lambda) = 1]\big| \leq \dfrac{q^2}{2} \cdot 2^{-n}$, which is negligible.
\end{framedlem}

\begin{proof}
  By the difference lemma it is enough to bound $\Pr[\mathrm{BAD}]$. The labels are $q$ independent uniform strings of $n$ bits. By the union bound over all pairs:
  \[
    \Pr[\mathrm{BAD}] = \Pr\big[\exists\, i \neq j : r_i = r_j\big]
    \leq \sum_{i \neq j} \Pr[r_i = r_j]
    = \binom{q}{2} \cdot \Col(\Uni{n})
    = \binom{q}{2} \cdot 2^{-n}
    \leq \frac{q^2}{2} \cdot 2^{-n}
  \]
  where $\Pr[r_i = r_j] = \Col(\Uni{n}) = 2^{-n}$: fixed $r_i$, the probability that $r_j$ is exactly the same string is one over $2^n$. Since $q = \poly(\lambda)$, $q^2$ is a polynomial and $2^{-n}$ is negligible, so the product is negligible. The bound holds even for unbounded adversaries: the two games are \textbf{statistically} close.
\end{proof}

\begin{framedobs}{It is the birthday paradox}
  With $23$ people there is already a $50\%$ chance that two share a birthday, because the \textbf{pairs} are many ($\binom{23}{2} = 253$). In the same way, collisions among the labels become likely when $q \approx 2^{n/2}$. With $n = 128$ that is about $2^{64}$ ciphertexts: in practice, this is the number of encryptions after which a key should be refreshed.
\end{framedobs}

\subsubsection{Step 3: putting the hops together}

\begin{proof}[Proof of \Cref{thm:prf-cpa}]
  By the triangle inequality, with \Cref{lem:game-hyb} and \Cref{lem:hyb-hybprime}:
  \[
    \begin{aligned}
      \big|\Pr[\mathrm{GAME}(\lambda, 0) = 1] - \Pr[\mathrm{GAME}(\lambda, 1) = 1]\big|
      &\leq \big|\Pr[\mathrm{GAME}(\lambda, 0) = 1] - \Pr[\mathrm{HYB}(\lambda, 0) = 1]\big| \\
      &\quad + \big|\Pr[\mathrm{HYB}(\lambda, 0) = 1] - \Pr[\mathrm{HYB}'(\lambda) = 1]\big| \\
      &\quad + \big|\Pr[\mathrm{HYB}'(\lambda) = 1] - \Pr[\mathrm{HYB}(\lambda, 1) = 1]\big| \\
      &\quad + \big|\Pr[\mathrm{HYB}(\lambda, 1) = 1] - \Pr[\mathrm{GAME}(\lambda, 1) = 1]\big| 
    \end{aligned}
  \]
  The first and the last terms are negligible by \Cref{lem:game-hyb}, the two middle terms by \Cref{lem:hyb-hybprime}. A sum of four negligible functions is negligible, so the scheme is CPA-secure.
\end{proof}

\begin{framedobs}{The structure of the proof}
  \begin{itemize}
    \item \textbf{Hop 1 (computational)} $\to$ replace the PRF by a truly random function; justified by the PRF assumption through a perfect reduction.
    \item \textbf{Hop 2 (statistical)} $\to$ replace every ciphertext by noise; no assumption at all, only the probability that two random labels collide.
    \item \textbf{The middle game does not depend on $b$} $\to$ it connects the game with $b = 0$ to the game with $b = 1$.
  \end{itemize}
\end{framedobs}

% ============================================================
%  The GGM construction
% ============================================================

\section{Pseudorandom functions from PRGs: the GGM construction}

\Cref{thm:prf-cpa} is useful only if PRFs exist. Goldreich, Goldwasser and Micali (GGM) showed how to build a PRF from any PRG, so PRFs exist whenever one-way functions do.

\subsection{Input of one bit}

Start from a \textbf{length-doubling} PRG and split its output into two halves of $\lambda$ bits:
\[
  G : \{0,1\}^\lambda \to \{0,1\}^{2\lambda} \qquad\qquad G(s) = \big(G_0(s),\ G_1(s)\big)
\]

\begin{framedprop}[label={prop:ggm-one-bit}]{A PRF with domain $\{0,1\}$}
  The family $F_k(x) = G_x(k)$, for $k \in \{0,1\}^\lambda$ and $x \in \{0,1\}$, is a PRF.
\end{framedprop}

The key is the seed of the PRG, and the input bit says \textbf{which half} to output: $F_k(0) = G_0(k)$ (left half) and $F_k(1) = G_1(k)$ (right half).

\textbf{Why it works.} A function with only two inputs is a table with \textbf{two rows}. A truly random such function is just two independent random values, that is, $2\lambda$ random bits. Here the whole table is $\big(F_k(0), F_k(1)\big) = G(k)$, which by the security of the PRG looks like $2\lambda$ random bits. So whatever the adversary asks, its view in $\mathrm{REAL}$ is indistinguishable from its view in $\mathrm{RAND}$.

\subsection{Input of $n$ bits: the tree}

One input bit is not enough. The idea of GGM is to repeat the choice ``left or right'' \textbf{once for every bit of the input}, building a binary tree:
\begin{itemize}
  \item the \textbf{root} is the key $k$;
  \item every node $s$ has two children: $G_0(s)$ on the left and $G_1(s)$ on the right;
  \item to compute $F_k(x_1 x_2 \cdots x_n)$, start from the root and go down: at level $i$ go \textbf{left if $x_i = 0$} and \textbf{right if $x_i = 1$};
  \item the output is the \textbf{leaf} reached.
\end{itemize}

\begin{frameddefn}[label={def:ggm}]{GGM pseudorandom function}
  Let $G(s) = (G_0(s), G_1(s))$ be a length-doubling PRG. For $k \in \{0,1\}^\lambda$ and $x = x_1 \cdots x_n \in \{0,1\}^n$:
  \[
    F_k(x_1 x_2 \cdots x_n) = G_{x_n}\Big(\cdots G_{x_2}\big(G_{x_1}(k)\big)\cdots\Big)
  \]
\end{frameddefn}

\begin{figure}[H]
  \centering
  \begin{tikzpicture}[
      level 1/.style={sibling distance=62mm, level distance=14mm},
      level 2/.style={sibling distance=31mm, level distance=14mm},
      every node/.style={font=\small},
      path/.style={very thick, BrickRed},
      edge from parent/.style={draw, thin}]
    \node {$k$}
      child {node {$G_0(k)$}
        child {node {$G_0(G_0(k))$}}
        child {node {$G_1(G_0(k))$}}
        edge from parent node[left, lbl] {$0$}}
      child {node[text=BrickRed] {$G_1(k)$}
        child {node[text=BrickRed] {$\mathbf{G_0(G_1(k))}$} edge from parent[path] node[left, lbl] {$0$}}
        child {node {$G_1(G_1(k))$}}
        edge from parent[path] node[right, lbl] {$1$}};
  \end{tikzpicture}
  \caption{The GGM tree for $n = 2$. The four leaves are the four values of $F_k$; the highlighted path computes $F_k(10) = G_0(G_1(k))$.}
\end{figure}

\begin{itemize}
  \item \textbf{The tree is the giant table} $\to$ it has $2^n$ leaves, one for every input.
  \item \textbf{It is efficient} $\to$ to compute $F_k(x)$ we only walk \textbf{one path}, so $n$ evaluations of $G$ are enough; the rest of the tree is never computed.
\end{itemize}

\begin{framedthm}[label={thm:ggm}]{Goldreich, Goldwasser, Micali}
  If $G$ is a length-doubling PRG, then the family of \Cref{def:ggm} is a PRF.
\end{framedthm}

\textbf{Proof idea.} A hybrid argument over the levels of the tree. In the hybrid $H_i$ the nodes of level $i$ are \textbf{truly random} (sampled lazily, only where the adversary's queries pass) and the levels below are computed with $G$. $H_0$ is the real PRF (only the root is random), and $H_n$ is a truly random function (all leaves are random). Two adjacent hybrids differ in replacing $G(s)$ by random strings on the nodes of one level that the adversary touches: these are at most polynomially many independent evaluations of $G$ on random seeds, which are indistinguishable from random by the security of the PRG (with a further hybrid over the evaluations). With $n$ levels and polynomially many queries, the total loss is negligible.

\begin{framedcor}{From one-way functions to CPA security}
  \[
    \text{OWF} \;\xrightarrow{\text{HILL}}\; \text{PRG} \;\xrightarrow{\text{GGM}}\; \text{PRF} \;\xrightarrow{\text{\Cref{thm:prf-cpa}}}\; \text{CPA-secure encryption}
  \]
  If one-way functions exist, then CPA-secure symmetric encryption exists.
\end{framedcor}

